\chapter{Similarity and Jordan Form}

In Chapters~26 and~27, we saw how cleanly symmetric and positive definite matrices behave. For symmetric matrices, the spectral theorem guarantees a full orthonormal basis of eigenvectors, real eigenvalues, and a diagonal factorization $\mat{A} = \mat{Q}\mat{\Lambda}\mat{Q}\T$; when the eigenvalues are positive, the associated quadratic forms curve upward in every direction like bowls.

Yet this chapter returns to a central question left open when we introduced diagonalization in Chapter~24: what happens when a square matrix cannot be diagonalized?

Diagonalization is the ideal outcome: in an eigenvector basis, an operator acts simply by scaling each coordinate axis independently. But not every matrix has a full basis of eigenvectors. The standard counterexample is a shear---a transformation that slides horizontal lines sideways proportionally to height:
\[
\mat{A} = \begin{bmatrix} 1 & 1 \\ 0 & 1 \end{bmatrix}.
\]
Its only eigenvalue is $1$ (with algebraic multiplicity $2$), but it possesses only \emph{one} line of eigenvectors, not two. Because there is no basis of eigenvectors, no diagonal matrix $\mat{D}$ exists. The transformation cannot be decomposed into independent stretches because its action includes a shear component along the invariant line.

This chapter investigates how to represent general linear operators as simply as possible. The underlying concept is \emph{similarity}: two matrices are similar when they represent the same linear operator under different coordinate systems. Diagonalization corresponds to the special case where the basis consists entirely of eigenvectors. When eigenvectors are missing, any square complex matrix can still be brought to a canonical, block-diagonal form called the \textbf{Jordan form}. It is diagonal except for superdiagonal $1$'s that couple generalized eigenvectors into chains. In this chapter, we develop the geometric meaning of similarity, the structure of generalized eigenvectors, and the classification of operators via Jordan blocks.

\section{Similarity: The Same Operator in a New Basis}
\label{ch28:similarity}

\subsection{The change-of-basis picture}

Here is the single idea the whole chapter rests on. A matrix is not really an operator;
it is the \emph{coordinates} of an operator once you have fixed a basis. Change the
basis and the same operator gets a different matrix. Two matrices that describe the
same underlying operator --- just in different coordinate systems --- are called
\emph{similar}.

\begin{intuitionbox}
Think of an operator $T$ as a physical motion of space: a rotation, a stretch, a shear.
The motion exists before you choose any axes. When you pick a basis and write down the
matrix, you are just recording that motion in a particular language. \textbf{Similar
matrices are translations of the same sentence into different languages.} The change of
basis is the dictionary.
\end{intuitionbox}

Here is the dictionary made explicit. Suppose $\mat{A}$ is the matrix of an operator in
the standard basis, and you want its matrix in a new basis whose vectors are the columns
of an invertible matrix $\mat{P}$. To apply the operator to a vector $\vec{y}$ given in
the new coordinates, you do three things:

\begin{enumerate}
\item \textbf{Translate into standard coordinates:} multiply by $\mat{P}$ to get
  $\mat{P}\vec{y}$, the same vector spelled in the standard basis.
\item \textbf{Apply the operator:} multiply by $\mat{A}$ to get $\mat{A}\mat{P}\vec{y}$.
\item \textbf{Translate the answer back into new coordinates:} multiply by $\mat{P}\inv$.
\end{enumerate}

The whole round trip is $\mat{P}\inv \mat{A}\, \mat{P}$, and that is the matrix of the
same operator in the new basis.

\begin{figure}[H]
\centering
\begin{tikzpicture}[>=Stealth, scale=1.0, every node/.style={font=\small}]
  \node[draw, rounded corners=2pt, fill=targetgreen!12, inner sep=6pt, align=center] (a) at (0,2.6)
    {New Coordinates\\$\vec{y}$};
  \node[draw, rounded corners=2pt, fill=targetgreen!12, inner sep=6pt, align=center] (b) at (6.5,2.6)
    {New Output\\$\mat{B}\vec{y}$};
  \node[draw, rounded corners=2pt, fill=sourceblue!12, inner sep=6pt, align=center] (c) at (0,0)
    {Standard Coordinates\\$\vec{x} = \mat{P}\vec{y}$};
  \node[draw, rounded corners=2pt, fill=sourceblue!12, inner sep=6pt, align=center] (d) at (6.5,0)
    {Standard Output\\$\mat{A}\vec{x} = \mat{A}\mat{P}\vec{y}$};

  \draw[->, very thick, structurepurple] (a) -- (b) node[midway, above=3pt, align=center, font=\footnotesize, text=structurepurple!90!black]
    {$\mat{B}=\mat{P}\inv\mat{A}\mat{P}$\\[1pt]\scriptsize (direct action)};
  \draw[->, thick, quotientorange] (a) -- (c) node[midway, left=3pt, font=\footnotesize] {$\mat{P}$ (change basis)};
  \draw[->, very thick, vecblue] (c) -- (d) node[midway, below=3pt, align=center, font=\footnotesize, text=vecblue!90!black]
    {$\mat{A}$\\[1pt]\scriptsize (standard action)};
  \draw[->, thick, quotientorange] (d) -- (b) node[midway, right=3pt, font=\footnotesize] {$\mat{P}\inv$ (return basis)};
\end{tikzpicture}
\caption{The change-of-basis commutative square. Applying $\mat{B}$ directly in the
new coordinates equals the three-step detour: translate to standard coordinates via $\mat{P}$,
apply the original transformation $\mat{A}$, then translate back via $\mat{P}\inv$. That
geometric equivalence is $\mat{B}=\mat{P}\inv\mat{A}\mat{P}$.}
\label{ch28:cob-square}
\end{figure}

\begin{definition}[Similar matrices]
\label{ch28:def-similar}
Two $n\times n$ matrices $\mat{A}$ and $\mat{B}$ are \textbf{similar}, written
$\mat{A}\sim\mat{B}$, if there is an invertible matrix $\mat{P}$ with
\[
\mat{B} = \mat{P}\inv \mat{A}\, \mat{P}.
\]
The matrix $\mat{P}$ is the change of basis: its columns are the new basis vectors
written in the old coordinates.
\end{definition}

Notice that diagonalization, $\mat{A}=\mat{P}\mat{D}\mat{P}\inv$, says exactly that
$\mat{D}=\mat{P}\inv\mat{A}\mat{P}$ --- that is, \emph{$\mat{A}$ is similar to a diagonal
matrix}. So ``diagonalizable'' is just a special, lucky case of similarity. The whole
chapter is the study of which canonical form you can reach when you are not so lucky.

\begin{focusbox}
$\mat{B}=\mat{P}\inv\mat{A}\mat{P}$ does not change \emph{what} the operator does. It
changes only the coordinate language used to describe it. Everything coordinate-free
about the operator --- its eigenvalues, its rank, its determinant, its trace --- must
therefore survive the change. That is the next thing we prove.
\end{focusbox}

\subsection{Similarity is an equivalence relation}

Before we mine similar matrices for shared properties, we should check that ``similar'' behaves
the way a notion of sameness should: every matrix is similar to itself, the relation is
symmetric, and it chains.

\begin{proposition}
\label{ch28:equiv}
Similarity is an equivalence relation: $\mat{A}\sim\mat{A}$; if $\mat{A}\sim\mat{B}$ then
$\mat{B}\sim\mat{A}$; and if $\mat{A}\sim\mat{B}$ and $\mat{B}\sim\mat{C}$ then
$\mat{A}\sim\mat{C}$.
\end{proposition}

\begin{proof}
\emph{Reflexive:} take $\mat{P}=\mat{I}$, so $\mat{A}=\mat{I}\inv\mat{A}\,\mat{I}$.
\emph{Symmetric:} if $\mat{B}=\mat{P}\inv\mat{A}\mat{P}$, multiply to solve for $\mat{A}$:
$\mat{A}=\mat{P}\mat{B}\mat{P}\inv=(\mat{P}\inv)\inv\mat{B}\,(\mat{P}\inv)$, so
$\mat{B}\sim\mat{A}$ via $\mat{P}\inv$. \emph{Transitive:} if
$\mat{B}=\mat{P}\inv\mat{A}\mat{P}$ and $\mat{C}=\mat{Q}\inv\mat{B}\mat{Q}$, substitute:
$\mat{C}=\mat{Q}\inv\mat{P}\inv\mat{A}\mat{P}\mat{Q}=(\mat{P}\mat{Q})\inv\mat{A}(\mat{P}\mat{Q})$,
so $\mat{A}\sim\mat{C}$ via $\mat{P}\mat{Q}$.
\end{proof}

So the set of all $n\times n$ matrices breaks into disjoint \emph{similarity classes},
each class being one operator seen in all possible bases. The grand goal of this chapter
is to pick, from each class, one cleanest representative. For nice operators that
representative is diagonal; in general it is the Jordan form.

\begin{recoverybox}
\textbf{Pause and think:} Suppose $\mat{A}$ and $\mat{B}$ are similar matrices. Could $\mat{A}$ be invertible while $\mat{B}$ is singular?
\emph{Reasoning:} No. Since $\mat{B}=\mat{P}\inv\mat{A}\mat{P}$, taking determinants immediately gives $\det\mat{B} = \det(\mat{P}\inv)\det(\mat{A})\det(\mat{P}) = \det\mat{A}$. If $\det\mat{A} \neq 0$, then $\det\mat{B} \neq 0$ as well. Even more directly, we can write down the inverse of $\mat{B}$ explicitly: $\mat{B}\inv = (\mat{P}\inv\mat{A}\mat{P})\inv = \mat{P}\inv\mat{A}\inv\mat{P}$. If one is invertible, both are invertible, and their inverses are themselves similar!
\end{recoverybox}

\section{What Similar Matrices Share}
\label{ch28:invariants}

\subsection{Eigenvalues, determinant, trace, rank}

Because $\mat{A}$ and $\mat{B}$ represent the same linear operator in different bases, anything that
is really an intrinsic property of the operator (and not of the chosen coordinates) must match. Here are four central invariants.

\begin{theorem}[Similarity invariants]
\label{ch28:thm-invariants}
If $\mat{B}=\mat{P}\inv\mat{A}\mat{P}$, then $\mat{A}$ and $\mat{B}$ have the same
characteristic polynomial $p(\lambda) = \det(\mat{A}-\lambda\mat{I})$ (the polynomial whose roots are the eigenvalues), hence the same eigenvalues with the same algebraic
multiplicities. They also have the same determinant, the same trace, and the same rank.
\end{theorem}

\begin{proof}
Compute the characteristic polynomial of $\mat{B}$ directly, using
$\mat{I}=\mat{P}\inv\mat{I}\mat{P}$:
\[
\det(\mat{B}-\lambda\mat{I})
= \det(\mat{P}\inv\mat{A}\mat{P}-\lambda\mat{P}\inv\mat{I}\mat{P})
= \det\!\big(\mat{P}\inv(\mat{A}-\lambda\mat{I})\mat{P}\big).
\]
Determinant is multiplicative, $\det(\mat{X}\mat{Y})=\det\mat{X}\det\mat{Y}$, and
$\det(\mat{P}\inv)=1/\det(\mat{P})$, so the two $\mat{P}$ factors cancel:
\[
\det(\mat{B}-\lambda\mat{I}) = \frac{1}{\det\mat{P}}\,\det(\mat{A}-\lambda\mat{I})\,\det\mat{P}
= \det(\mat{A}-\lambda\mat{I}).
\]
Same characteristic polynomial. Its roots (eigenvalues) and their multiplicities are
therefore identical. Setting $\lambda=0$ gives $\det\mat{B}=\det\mat{A}$. The trace
matches because $\tr(\mat{X}\mat{Y})=\tr(\mat{Y}\mat{X})$ for any conformable
$\mat{X},\mat{Y}$, so
$\tr(\mat{P}\inv\mat{A}\mat{P})=\tr(\mat{A}\mat{P}\mat{P}\inv)=\tr\mat{A}$. Finally,
multiplying by the invertible matrices $\mat{P}$ and $\mat{P}\inv$ never changes rank, so
$\rank\mat{B}=\rank\mat{A}$.
\end{proof}

\begin{labbox}
Verify the trace and determinant are preserved with your hands. Take
$\mat{A}=\begin{bmatrix}2&3\\0&3\end{bmatrix}$ and the shear
$\mat{P}=\begin{bmatrix}1&1\\0&1\end{bmatrix}$ (so $\mat{P}\inv=\begin{bmatrix}1&-1\\0&1\end{bmatrix}$).
Compute $\mat{B}=\mat{P}\inv\mat{A}\mat{P}$ by hand, in two steps:
$\mat{A}\mat{P}=\begin{bmatrix}2&5\\0&3\end{bmatrix}$, then
$\mat{P}\inv(\mat{A}\mat{P})=\begin{bmatrix}2&2\\0&3\end{bmatrix}=\mat{B}$.
Check: $\tr=5$ and $\det=6$ for both, and both have eigenvalues $2,3$. The off-diagonal
entry changed from $3$ to $2$; the operator did not. (Watch the order: $\mat{B}$ is
\emph{not} the same as $\mat{A}$ here, only similar to it.)
\end{labbox}

\begin{trapbox}
Sharing eigenvalues is \emph{necessary} for similarity but \textbf{not sufficient}. The
two matrices
\[
\begin{bmatrix}1&0\\0&1\end{bmatrix}
\qquad\text{and}\qquad
\begin{bmatrix}1&1\\0&1\end{bmatrix}
\]
both have the single eigenvalue $1$ with algebraic multiplicity $2$, the same trace, the
same determinant, the same characteristic polynomial. Yet they are \emph{not} similar:
the first is the identity (it commutes with everything, so
$\mat{P}\inv\mat{I}\mat{P}=\mat{I}$ always), while the second is a nontrivial shear.
Eigenvalues alone do not pin down a similarity class. We will need more --- the Jordan
form --- to tell these apart.
\end{trapbox}

\subsection{Eigenvectors transform too}

Eigenvalues are shared, but the eigenvectors get \emph{relabeled} by the change of
basis, which makes perfect sense: an eigenvector is a direction in space, and you wrote
it in new coordinates.

\begin{proposition}
\label{ch28:eigvec-transform}
If $\mat{A}\vec{v}=\lambda\vec{v}$ and $\mat{B}=\mat{P}\inv\mat{A}\mat{P}$, then
$\vec{w}=\mat{P}\inv\vec{v}$ is an eigenvector of $\mat{B}$ with the same eigenvalue
$\lambda$.
\end{proposition}

\begin{proof}
$\mat{B}\vec{w}=\mat{P}\inv\mat{A}\mat{P}\,(\mat{P}\inv\vec{v})
=\mat{P}\inv\mat{A}\vec{v}=\mat{P}\inv(\lambda\vec{v})=\lambda(\mat{P}\inv\vec{v})=\lambda\vec{w}$.
\end{proof}

So eigenvalues live in the operator and survive any change of basis; eigenvectors live
in the operator too, but their \emph{coordinates} change with the basis. This is exactly
why diagonalization works: the special basis is the eigenvector basis, and in those
coordinates each eigenvector is a standard axis $\vec{e}_i$, on which the operator acts
by simple scaling.

\begin{strangbox}
Strang introduces $\mat{A}=\mat{S}\mat{\Lambda}\mat{S}\inv$ (he writes $\mat{S}$ for the
eigenvector matrix and $\mat{\Lambda}$ for the diagonal of eigenvalues). Read that
equation as a similarity statement: $\mat{\Lambda}=\mat{S}\inv\mat{A}\mat{S}$, so
$\mat{A}$ \emph{is similar to} $\mat{\Lambda}$. The present chapter is what happens when
$\mat{S}$ cannot be formed because there are too few eigenvectors --- the case Strang
flags as ``not diagonalizable'' and largely sets aside. Here we follow it to its
conclusion.
\end{strangbox}

\section{When Diagonalization Fails}
\label{ch28:fails}

\subsection{Algebraic versus geometric multiplicity}

A matrix is diagonalizable exactly when it has a full basis of eigenvectors --- $n$
independent ones for an $n\times n$ matrix. The way that can fail is subtle, and pinning
it down means separating two different counts attached to each eigenvalue $\lambda$.

\begin{definition}[Algebraic and geometric multiplicity]
\label{ch28:def-mult}
Let $\lambda$ be an eigenvalue of $\mat{A}$.
\begin{itemize}
\item Its \textbf{algebraic multiplicity} is the number of times $(\,t-\lambda\,)$ divides
  the characteristic polynomial $\det(\mat{A}-t\mat{I})$ --- how many times $\lambda$
  shows up as a root.
\item Its \textbf{geometric multiplicity} is $\dimn\nullsp(\mat{A}-\lambda\mat{I})$ ---
  the dimension of the eigenspace, i.e.\ how many independent eigenvectors $\lambda$
  actually supplies.
\end{itemize}
\end{definition}

It is a theorem (we will not belabor it) that for every eigenvalue,
\[
1 \;\le\; \text{geometric multiplicity} \;\le\; \text{algebraic multiplicity}.
\]
You always have at least one eigenvector, and you never have more independent
eigenvectors than the root multiplicity allows. \textbf{Diagonalization succeeds exactly
when geometric equals algebraic for every eigenvalue.} When some eigenvalue is
``geometrically deficient'' --- fewer eigenvectors than its algebraic multiplicity
demands --- the eigenvector basis is short, and no diagonal form exists. That gap
(algebraic multiplicity minus geometric multiplicity) is called the \emph{defect} of the eigenvalue.

\begin{intuitionbox}
The algebraic multiplicity is the multiplicity of $\lambda$ as a root of the
characteristic polynomial. The geometric multiplicity counts how many independent
eigenvector directions the matrix provides. When geometric multiplicity is strictly less
than algebraic multiplicity, the matrix is defective, and the deficiency reflects shearing behavior.
\end{intuitionbox}

\subsection{The canonical defective matrix}

Look hard at the smallest defective matrix, the shear we met at the start:
\[
\mat{J} = \begin{bmatrix} 1 & 1 \\ 0 & 1 \end{bmatrix}.
\]
Its characteristic polynomial is $\det\!\begin{bmatrix}1-t&1\\0&1-t\end{bmatrix}=(1-t)^2$,
so $\lambda=1$ has algebraic multiplicity $2$. Now hunt for eigenvectors. Solve
$(\mat{J}-\mat{I})\vec{v}=\vec{0}$:
\[
\mat{J}-\mat{I}=\begin{bmatrix}0&1\\0&0\end{bmatrix},\qquad
\begin{bmatrix}0&1\\0&0\end{bmatrix}\colvec{x\\y}=\colvec{y\\0}=\colvec{0\\0}
\;\Longrightarrow\; y=0.
\]
Every eigenvector is a multiple of $\colvec{1\\0}$. The eigenspace is one-dimensional:
geometric multiplicity $1$, algebraic multiplicity $2$. Defect $=1$. There is no second
independent eigenvector anywhere, so no eigenvector basis, so no diagonalization. The
matrix is genuinely, irreparably non-diagonal.

\begin{figure}[H]
\centering
\begin{tikzpicture}[>=Stealth,scale=1.1,every node/.style={font=\small}]
\begin{scope}[xshift=0cm]
  \draw[step=0.5,gray!25,very thin] (-0.2,-0.2) grid (2.4,1.8);
  \draw[->,gray!60] (-0.3,0)--(2.6,0) node[right]{$x$};
  \draw[->,gray!60] (0,-0.3)--(0,2.0) node[above]{$y$};
  \fill[sourceblue!12] (0,0) -- (1,0) -- (1,1) -- (0,1) -- cycle;
  \draw[sourceblue,thick] (0,0) -- (1,0) -- (1,1) -- (0,1) -- cycle;
  \draw[->,very thick,targetgreen] (0,0)--(1,0) node[midway,below=2pt]{$\vec{e}_1$};
  \draw[->,very thick,structurepurple] (0,0)--(0,1) node[midway,left=2pt]{$\vec{e}_2$};
  \node at (1.1,-0.75) {before: unit square};
\end{scope}
\begin{scope}[xshift=5.2cm]
  \draw[step=0.5,gray!20,very thin,dashed] (-0.2,-0.2) grid (2.6,1.8);
  \draw[->,gray!60] (-0.3,0)--(2.8,0) node[right]{$x$};
  \draw[->,gray!60] (0,-0.3)--(0,2.0) node[above]{$y$};
  \fill[vecred!12] (0,0) -- (1,0) -- (2,1) -- (1,1) -- cycle;
  \draw[vecred,thick] (0,0) -- (1,0) -- (2,1) -- (1,1) -- cycle;
  \draw[->,very thick,targetgreen] (0,0)--(1,0) node[midway,below=2pt]{$\mat{J}\vec{e}_1=\vec{e}_1$};
  \draw[->,very thick,structurepurple] (0,0)--(1,1) node[midway,above left=1pt,fill=white,inner sep=1pt]{$\mat{J}\vec{e}_2$};
  \draw[->,thick,quotientorange,dashed] (0,1)--(1,1) node[midway,above=1pt,font=\footnotesize]{slide by $y=1$};
  \node at (1.3,-0.75) {after shear: slides right};
\end{scope}
\end{tikzpicture}
\caption{The shear $\mat{J}=\protect\begin{bmatrix}1&1\\0&1\protect\end{bmatrix}$ slides
each point $(x,y)$ to $(x+y,y)$. The horizontal axis stays fixed: $\vec{e}_1=\protect\colvec{1\\0}$
is the sole eigenvector direction. The vertical vector $\vec{e}_2=\protect\colvec{0\\1}$ is tilted
to $\mat{J}\vec{e}_2 = \protect\colvec{1\\1} = \vec{e}_1 + \vec{e}_2$. Because every direction
off the horizontal axis is tilted toward the horizontal, no second independent eigenvector can exist.}
\label{ch28:shear-grid}
\end{figure}

\begin{walkbox}
Take a walk and ask: why \emph{should} a shear have only one eigenvector? An eigenvector
is a direction that the operator merely scales --- it does not get rotated off its own
line. A shear pushes every direction except the horizontal a little toward the
horizontal. So no slanted direction can survive unrotated, and only the horizontal axis
is left. The geometry forces the deficiency; the algebra merely records it.
\end{walkbox}

\section{Generalized Eigenvectors: Filling the Gap}
\label{ch28:generalized}

Because $\mat{J}$ has only one independent eigenvector, we broaden the definition to vectors
annihilated by higher powers of $\mat{J}-\lambda\mat{I}$.
A true eigenvector satisfies $(\mat{J}-\mat{I})\vec{v}=\vec{0}$.
A \emph{generalized} eigenvector satisfies
$(\mat{J}-\mat{I})^k\vec{v}=\vec{0}$ for some integer $k\ge 1$, even if
$(\mat{J}-\mat{I})\vec{v}\neq\vec{0}$.

\begin{definition}[Generalized eigenvector]
\label{ch28:def-gen}
A nonzero vector $\vec{v}$ is a \textbf{generalized eigenvector} of $\mat{A}$ for the
eigenvalue $\lambda$ if $(\mat{A}-\lambda\mat{I})^k\vec{v}=\vec{0}$ for some integer
$k\ge 1$. The smallest such $k$ is its \emph{rank} (or \emph{height}). An ordinary
eigenvector is the case $k=1$.
\end{definition}

To find one for $\mat{J}$, we need a vector that $(\mat{J}-\mat{I})$ does \emph{not}
kill but $(\mat{J}-\mat{I})^2$ does. First note
\[
(\mat{J}-\mat{I})^2=\begin{bmatrix}0&1\\0&0\end{bmatrix}\begin{bmatrix}0&1\\0&0\end{bmatrix}
=\begin{bmatrix}0&0\\0&0\end{bmatrix},
\]
the zero matrix. So $(\mat{J}-\mat{I})^2$ kills \emph{everything} --- every vector in
$\R^2$ is a generalized eigenvector of height at most $2$. Pick one that the first power
does not kill, for instance $\vec{w}=\colvec{0\\1}$:
\[
(\mat{J}-\mat{I})\vec{w}=\begin{bmatrix}0&1\\0&0\end{bmatrix}\colvec{0\\1}=\colvec{1\\0}=\vec{v}\neq\vec{0}.
\]
So the first power of $(\mat{J}-\mat{I})$ does not annihilate $\vec{w}$; instead it
sends $\vec{w}$ onto the genuine eigenvector $\vec{v}=\colvec{1\\0}$. The eigenvector and
the generalized eigenvector form a little chain
\[
\vec{w}\;\xrightarrow{\;\mat{J}-\mat{I}\;}\;\vec{v}\;\xrightarrow{\;\mat{J}-\mat{I}\;}\;\vec{0}.
\]
This chain is the missing structure. The pair $\{\vec{v},\vec{w}\}$ is independent (one
is $\colvec{1\\0}$, the other $\colvec{0\\1}$), so it \emph{is} a basis of $\R^2$ --- not
a basis of eigenvectors, but a basis of generalized eigenvectors. That is the
replacement for the eigenvector basis we could not build.

\begin{intuitionbox}
A generalized eigenvector is a vector that is ``almost'' an eigenvector --- it is not
fixed by the operator, but applying $(\mat{A}-\lambda\mat{I})$ repeatedly walks it down a
staircase of vectors until it lands on a true eigenvector and then on zero. The length
of that staircase is exactly the size of the Jordan block, and the staircase is called a
\textbf{Jordan chain}.
\end{intuitionbox}

\subsection{What the chain does to the matrix}

Now build the change of basis from this chain and see what falls out. Order the chain
``bottom up'': put the eigenvector first, the generalized eigenvector second. So the new
basis is $\vec{v}=\colvec{1\\0}$ then $\vec{w}=\colvec{0\\1}$, giving
\[
\mat{P}=\begin{bmatrix}\vec{v}&\vec{w}\end{bmatrix}=\begin{bmatrix}1&0\\0&1\end{bmatrix}=\mat{I}.
\]
(For this particular $\mat{J}$ the chain vectors happen to be the standard basis already,
which is why $\mat{J}$ is its own Jordan form.) The defining relations of the chain are
\[
\mat{J}\vec{v}=\vec{v},\qquad \mat{J}\vec{w}=\vec{v}+\vec{w}.
\]
Read those as columns: the operator sends the first basis vector to $1\cdot\vec{v}$, and
the second basis vector to $1\cdot\vec{v}+1\cdot\vec{w}$. Writing those coefficients as
columns gives the matrix in the chain basis:
\[
\begin{bmatrix}1&1\\0&1\end{bmatrix}.
\]
The eigenvalue $\lambda=1$ sits on the diagonal, and the single $1$ \emph{above} the
diagonal is the algebraic shadow of the relation $\mat{J}\vec{w}=\lambda\vec{w}+\vec{v}$
--- the ``$+\vec{v}$'' that says $\vec{w}$ is dragged partly onto the eigenvector. That
extra $1$ is the slide. It is what diagonalization could never capture, and the Jordan
form keeps it on purpose.

\section{The Jordan Form}
\label{ch28:jordan}

\subsection{Jordan blocks}

The atom of the theory is a single Jordan block: an eigenvalue $\lambda$ down the
diagonal, with $1$'s on the superdiagonal.

\begin{definition}[Jordan block]
\label{ch28:def-block}
A \textbf{Jordan block} $\mat{J}_k(\lambda)$ is the $k\times k$ matrix with $\lambda$ on
the diagonal, $1$ on the entries just above the diagonal, and $0$ everywhere else:
\[
\mat{J}_1(\lambda)=\begin{bmatrix}\lambda\end{bmatrix},\quad
\mat{J}_2(\lambda)=\begin{bmatrix}\lambda&1\\0&\lambda\end{bmatrix},\quad
\mat{J}_3(\lambda)=\begin{bmatrix}\lambda&1&0\\0&\lambda&1\\0&0&\lambda\end{bmatrix},\ \dots
\]
\end{definition}

A single Jordan block has exactly \emph{one} eigenvector (geometric multiplicity $1$) no
matter how big it is, but algebraic multiplicity $k$. So a block of size $k$ carries a
defect of $k-1$: one eigenvector and $k-1$ generalized eigenvectors, strung in one chain
of length $k$.

\begin{labbox}
Confirm the eigenvector count for $\mat{J}_3(5)$. Compute
$\mat{J}_3(5)-5\mat{I}=\begin{bmatrix}0&1&0\\0&0&1\\0&0&0\end{bmatrix}$ and solve
$(\mat{J}_3(5)-5\mat{I})\vec{x}=\vec{0}$. You should find the only constraints are
$x_2=0$ and $x_3=0$, so the eigenspace is the single line $\spn\{\vec{e}_1\}$ ---
geometric multiplicity $1$, even though $5$ is a triple root. Then check that
$(\mat{J}_3(5)-5\mat{I})^3=\mat{0}$ but the square is not zero: the chain has length
exactly $3$.
\end{labbox}

\subsection{The big theorem}

Every operator, over the complex numbers, can be assembled out of these blocks --- one
diagonal staircase per Jordan chain --- and that assembly is essentially unique.

\begin{theorem}[Jordan canonical form]
\label{ch28:thm-jordan}
Let $\mat{A}$ be any $n\times n$ matrix over $\C$. Then $\mat{A}$ is similar to a
block-diagonal matrix
\[
\mat{J} = \begin{bmatrix}
\mat{J}_{k_1}(\lambda_1) & & \\
 & \ddots & \\
 & & \mat{J}_{k_m}(\lambda_m)
\end{bmatrix},
\]
a direct sum of Jordan blocks. The list of blocks --- their eigenvalues and sizes --- is
unique up to the order in which the blocks are written. We call $\mat{J}$ the
\textbf{Jordan canonical form} of $\mat{A}$, and $\mat{A}=\mat{P}\mat{J}\mat{P}\inv$ for
some invertible $\mat{P}$ whose columns are the Jordan chains.
\end{theorem}

\begin{figure}[H]
\centering
\begin{tikzpicture}[scale=0.9, every node/.style={font=\small}]
  \node[font=\normalsize\bfseries] at (-0.7, 2.5) {$\mat{J} = $};
  \draw[thick, spacegray] (0.15,0.05) -- (0,0.05) -- (0,4.95) -- (0.15,4.95);
  \draw[thick, spacegray] (7.85,0.05) -- (8.0,0.05) -- (8.0,4.95) -- (7.85,4.95);

  \draw[fill=sourceblue!14, draw=sourceblue, very thick, rounded corners=3pt] (0.3, 2.2) rectangle (3.5, 4.8);
  \node[sourceblue!80!black, font=\footnotesize\bfseries] at (1.9, 4.45) {$\mat{J}_3(\lambda_1)$};
  \node at (0.9, 3.8) {$\lambda_1$};
  \node[vecred, font=\footnotesize\bfseries] at (1.9, 3.8) {$1$};
  \node at (2.9, 3.8) {$0$};
  \node at (0.9, 3.1) {$0$};
  \node at (1.9, 3.1) {$\lambda_1$};
  \node[vecred, font=\footnotesize\bfseries] at (2.9, 3.1) {$1$};
  \node at (0.9, 2.4) {$0$};
  \node at (1.9, 2.4) {$0$};
  \node at (2.9, 2.4) {$\lambda_1$};

  \draw[fill=targetgreen!14, draw=targetgreen, very thick, rounded corners=3pt] (3.8, 0.8) rectangle (6.2, 2.1);
  \node[targetgreen!80!black, font=\footnotesize\bfseries] at (5.0, 1.85) {$\mat{J}_2(\lambda_2)$};
  \node at (4.4, 1.5) {$\lambda_2$};
  \node[vecred, font=\footnotesize\bfseries] at (5.6, 1.5) {$1$};
  \node at (4.4, 1.05) {$0$};
  \node at (5.6, 1.05) {$\lambda_2$};

  \draw[fill=quotientorange!14, draw=quotientorange, very thick, rounded corners=3pt] (6.5, 0.15) rectangle (7.7, 0.75);
  \node[quotientorange!80!black, font=\footnotesize\bfseries] at (7.1, 0.45) {$[\lambda_3]$};

  \node[gray!50, font=\Large\bfseries] at (5.6, 3.5) {$\mat{0}$};
  \node[gray!50, font=\Large\bfseries] at (2.0, 1.1) {$\mat{0}$};

  \node[anchor=west, align=left, font=\footnotesize] at (8.5, 4.0)
    {\textbf{Diagonal blocks:}\\independent chains\\for eigenvalues $\lambda_i$};
  \node[anchor=west, align=left, font=\footnotesize] at (8.5, 2.5)
    {\textbf{Superdiagonal $1$s:}\\shearing slide coupling\\chain vectors};
  \node[anchor=west, align=left, font=\footnotesize] at (8.5, 1.0)
    {\textbf{Off-block zeros:}\\no interaction between\\different blocks};
\end{tikzpicture}
\caption{The anatomy of a matrix in Jordan canonical form. The operator decomposes into
independent diagonal blocks $\mat{J}_k(\lambda)$. Within each block, the eigenvalue
$\lambda$ acts by pure scaling along the diagonal, while the superdiagonal $1$'s couple
successive vectors along a Jordan chain. All off-block entries are strictly zero.}
\label{ch28:jordan-anatomy}
\end{figure}

We will not prove existence in full --- the proof is a careful induction on generalized
eigenspaces, more machinery than a first reader needs. But every \emph{idea} in the proof
is already in front of you: find the eigenvectors; where you come up short, climb to
generalized eigenvectors; bundle each chain into a block. What you should trust and use is the
bookkeeping, which is entirely concrete.

\begin{formalbox}
The block structure is read off from ranks of powers of $(\mat{A}-\lambda\mat{I})$. For
a fixed eigenvalue $\lambda$:
\begin{itemize}
\item The \textbf{number of Jordan blocks} for $\lambda$ equals the geometric
  multiplicity $\dimn\nullsp(\mat{A}-\lambda\mat{I})$ --- one block per genuine
  eigenvector.
\item The \textbf{total size} of all blocks for $\lambda$ equals its algebraic
  multiplicity.
\item The \textbf{size of the largest block} for $\lambda$ equals the smallest $k$ with
  $\nullsp(\mat{A}-\lambda\mat{I})^k=\nullsp(\mat{A}-\lambda\mat{I})^{k+1}$ --- the point
  where the null spaces stop growing.
\end{itemize}
The exact sizes come from the chain of nullity dimensions
$d_j=\dimn\nullsp(\mat{A}-\lambda\mat{I})^j$. The number of blocks of size $\ge j$ is
$d_j-d_{j-1}$, so the number of blocks of size exactly $j$ is
$(d_j-d_{j-1})-(d_{j+1}-d_j) = 2d_j - d_{j-1} - d_{j+1}$.
\end{formalbox}

\begin{focusbox}
Two numbers per eigenvalue carry most of the meaning. Algebraic multiplicity =
\emph{total} room the eigenvalue occupies on the diagonal. Geometric multiplicity =
\emph{number of blocks}. Diagonalizable means every block has size $1$, which happens
exactly when these two numbers are equal for every eigenvalue. The further apart they
are, the bigger the blocks, the more shearing the operator hides.
\end{focusbox}

A picture makes the bookkeeping stick. Stack the blocks for a fixed
eigenvalue as rows of boxes --- one row per block, the row's length equal to that block's
size. The \emph{number of rows} is the geometric multiplicity (one row per genuine
eigenvector); the \emph{total number of boxes} is the algebraic multiplicity. The growing
nullities $d_j$ are then just the heights of the columns: $d_1$ is the number of rows
(boxes in the first column), $d_2-d_1$ is the boxes added in the second column, and so on.
Reading the diagram by rows gives block sizes; reading it by columns gives the nullity
jumps. The same staircase, two ways.

\begin{figure}[H]
\centering
\begin{tikzpicture}[scale=0.85,every node/.style={font=\small}]
% A partition 3+1 of algebraic multiplicity 4 (two blocks)
\foreach \x in {0,1,2}{ \draw[fill=structurepurple!15,draw=structurepurple!80,thick] (\x,1) rectangle ++(1,1); }
\foreach \x in {0}{ \draw[fill=structurepurple!15,draw=structurepurple!80,thick] (\x,0) rectangle ++(1,1); }
\node at (1.5,2.45) {block sizes $3+1$};
\node[anchor=east] at (-0.2,1.5) {block 1};
\node[anchor=east] at (-0.2,0.5) {block 2};
% column annotations
\draw[->,vecblue,thick] (0.5,-0.5)--(0.5,-0.05);
\node[vecblue,font=\footnotesize] at (0.5,-0.85) {$d_1=2$};
\draw[->,targetgreen,thick] (1.5,-0.5)--(1.5,0.95);
\node[targetgreen,font=\footnotesize] at (1.5,-0.85) {$+1$};
\draw[->,quotientorange,thick] (2.5,-0.5)--(2.5,0.95);
\node[quotientorange,font=\footnotesize] at (2.5,-0.85) {$+1$};
\node[anchor=west,align=left] at (3.8,1.0)
  {rows $=$ geometric mult.\ $=2$\\[3pt]
   boxes $=$ algebraic mult.\ $=4$\\[3pt]
   column heights $=d_1,\,d_2{-}d_1,\dots$};
\end{tikzpicture}
\caption{The Jordan structure for one eigenvalue drawn as a partition. Each row is a block;
its length is the block size. Rows count genuine eigenvectors (geometric multiplicity);
total boxes count algebraic multiplicity. The column heights are exactly the nullity
increments $d_1,\,d_2-d_1,\,d_3-d_2,\dots$, which is why the formula
$2d_j-d_{j-1}-d_{j+1}$ counts blocks of a given size.}
\label{ch28:partition}
\end{figure}

\begin{recoverybox}
\textbf{Pause and think:} A $4\times 4$ matrix has a single eigenvalue $\lambda=3$ of algebraic multiplicity $4$. You compute the rank of $\mat{A}-3\mat{I}$ and find it is $2$. What are the possible Jordan forms?
\emph{Reasoning:} By the rank-nullity theorem, the geometric multiplicity is $\dimn\nullsp(\mat{A}-3\mat{I}) = 4 - 2 = 2$, which means there are exactly $2$ Jordan blocks. The block sizes must sum to the algebraic multiplicity $4$. The only partitions of $4$ into two positive integers are $3+1$ and $2+2$. Thus, the Jordan form is either $\mat{J}_3(3)\oplus\mat{J}_1(3)$ or $\mat{J}_2(3)\oplus\mat{J}_2(3)$. One more rank calculation---evaluating the nullity of $(\mat{A}-3\mat{I})^2$---decides between them immediately: if $d_2 = 3$, the partition is $3+1$; if $d_2 = 4$, it is $2+2$.
\end{recoverybox}

\subsection{A worked $3\times 3$ example}

Run the whole machine on a concrete matrix that is \emph{not} diagonalizable. Take
\[
\mat{A}=\begin{bmatrix}5&1&0\\0&5&0\\0&0&5\end{bmatrix}.
\]
The only eigenvalue is $5$, with algebraic multiplicity $3$ (it is upper triangular, so
read the diagonal). Now the eigenvectors. Form
\[
\mat{N}=\mat{A}-5\mat{I}=\begin{bmatrix}0&1&0\\0&0&0\\0&0&0\end{bmatrix}.
\]
Solving $\mat{N}\vec{x}=\vec{0}$ forces $x_2=0$; $x_1$ and $x_3$ are free. So the
eigenspace is two-dimensional, $\spn\{\vec{e}_1,\vec{e}_3\}$ --- geometric multiplicity
$2$. By the bookkeeping rules: \textbf{2 blocks} (geometric multiplicity), \textbf{total
size 3} (algebraic multiplicity). The only way to split $3$ into two positive block
sizes is $2+1$. So the Jordan form is
\[
\mat{J}=\begin{bmatrix}5&1&0\\0&5&0\\0&0&5\end{bmatrix}
=\mat{J}_2(5)\oplus\mat{J}_1(5).
\]
Here $\mat{A}$ already \emph{is} in Jordan form. Identifying the chains shows you
where the columns of $\mat{P}$ come from. We need one chain of length $2$ and one of
length $1$.

\emph{Length-1 chain (the lone eigenvector):} $\vec{e}_3=\colvec{0\\0\\1}$ satisfies
$\mat{N}\vec{e}_3=\vec{0}$. That is a block $\mat{J}_1(5)$ all by itself.

\emph{Length-2 chain:} we need a generalized eigenvector $\vec{w}$ with
$\mat{N}\vec{w}\neq\vec{0}$ but $\mat{N}^2\vec{w}=\vec{0}$. Since $\mat{N}^2=\mat{0}$
(check it), every vector dies by the second power; we just need one the first power does
not kill. Try $\vec{w}=\vec{e}_2=\colvec{0\\1\\0}$:
\[
\mat{N}\vec{e}_2=\colvec{1\\0\\0}=\vec{e}_1\neq\vec{0}.
\]
So the chain is $\vec{e}_2\to\vec{e}_1\to\vec{0}$, with eigenvector $\vec{e}_1$ at the
bottom. Assembling the basis (each chain bottom-up, blocks side by side) gives
$\mat{P}=[\,\vec{e}_1\ \vec{e}_2\ \vec{e}_3\,]=\mat{I}$, confirming $\mat{A}=\mat{J}$.

\begin{example}[A defective matrix that is \emph{not} already in Jordan form]
\label{ch28:ex-notjordan}
Now a less tidy one, where $\mat{P}\neq\mat{I}$:
\[
\mat{A}=\begin{bmatrix}3&1\\-1&5\end{bmatrix}.
\]
Characteristic polynomial: $\det\!\begin{bmatrix}3-t&1\\-1&5-t\end{bmatrix}
=(3-t)(5-t)+1=t^2-8t+16=(t-4)^2$. Eigenvalue $4$, algebraic multiplicity $2$. Eigenvectors:
\[
\mat{N}=\mat{A}-4\mat{I}=\begin{bmatrix}-1&1\\-1&1\end{bmatrix},\qquad
\mat{N}\vec{x}=\vec{0}\;\Longrightarrow\; -x_1+x_2=0\;\Longrightarrow\; x_2=x_1.
\]
The eigenspace is the single line $\spn\{\colvec{1\\1}\}$ --- geometric multiplicity $1$,
so \textbf{one block} of \textbf{size 2}: $\mat{J}=\mat{J}_2(4)=\begin{bmatrix}4&1\\0&4\end{bmatrix}$.
Build the chain. The eigenvector is $\vec{v}=\colvec{1\\1}$. For the generalized
eigenvector we need $\vec{w}$ with $\mat{N}\vec{w}=\vec{v}$. Solve
\[
\begin{bmatrix}-1&1\\-1&1\end{bmatrix}\colvec{w_1\\w_2}=\colvec{1\\1}
\;\Longrightarrow\; -w_1+w_2=1.
\]
Any solution works; pick $\vec{w}=\colvec{0\\1}$. Then
\[
\mat{P}=\begin{bmatrix}\vec{v}&\vec{w}\end{bmatrix}=\begin{bmatrix}1&0\\1&1\end{bmatrix},
\qquad \mat{P}\inv=\begin{bmatrix}1&0\\-1&1\end{bmatrix}.
\]
Verify the similarity:
\[
\mat{P}\inv\mat{A}\mat{P}
=\begin{bmatrix}1&0\\-1&1\end{bmatrix}\begin{bmatrix}3&1\\-1&5\end{bmatrix}\begin{bmatrix}1&0\\1&1\end{bmatrix}
=\begin{bmatrix}1&0\\-1&1\end{bmatrix}\begin{bmatrix}4&1\\4&5\end{bmatrix}
=\begin{bmatrix}4&1\\0&4\end{bmatrix}=\mat{J}_2(4).\ \checkmark
\]
The off-diagonal $1$ is exactly the relation $\mat{A}\vec{w}=4\vec{w}+\vec{v}$ written in
the chain basis.
\end{example}

Figure~\ref{ch28:concrete-chain} shows this chain living in the plane. The eigenvector
$\vec{v}=\colvec{1\\1}$ marks the one fixed direction; the generalized eigenvector
$\vec{w}=\colvec{0\\1}$ is off that line, and $\mat{N}=\mat{A}-4\mat{I}$ pushes it exactly
back onto $\vec{v}$.

\begin{figure}[H]
\centering
\begin{tikzpicture}[>=Stealth,scale=1.3,every node/.style={font=\small}]
\draw[->,gray!60] (-0.3,0)--(2.0,0) node[right]{$x_1$};
\draw[->,gray!60] (0,-0.3)--(0,2.0) node[above]{$x_2$};
\draw[vecblue,dashed] (-0.2,-0.2)--(1.8,1.8);
\node[vecblue, right] at (1.75,1.75) {eigenspace $x_2=x_1$};
\draw[->,very thick,vecblue] (0,0)--(1,1) node[below right]{$\vec{v}$};
\draw[->,very thick,structurepurple] (0,0)--(0,1) node[above left]{$\vec{w}$};
\draw[->,thick,targetgreen,dotted] (0,1) to[bend left=20] node[above,pos=0.6]{$\mat{N}\vec{w}=\vec{v}$} (1,1);
\end{tikzpicture}
\caption{The Jordan chain of Example~\ref{ch28:ex-notjordan} drawn in $\R^2$. The
eigenvector $\vec{v}=\protect\colvec{1\\1}$ (blue) spans the only fixed line. The
generalized eigenvector $\vec{w}=\protect\colvec{0\\1}$ (purple) sits off that line, and
$\mat{N}=\mat{A}-4\mat{I}$ carries $\vec{w}$ straight onto $\vec{v}$ (green) --- one step
down the staircase before the eigenvector is annihilated.}
\label{ch28:concrete-chain}
\end{figure}

\section{The Geometry of a Jordan Block}
\label{ch28:geometry}

A diagonal entry $\lambda$ scales an axis. The superdiagonal $1$ adds a \emph{nilpotent
shift}. Splitting a $2\times 2$ block displays this separation directly:
\[
\mat{J}_2(\lambda)=\begin{bmatrix}\lambda&1\\0&\lambda\end{bmatrix}
=\underbrace{\begin{bmatrix}\lambda&0\\0&\lambda\end{bmatrix}}_{\text{scale by }\lambda}
+\underbrace{\begin{bmatrix}0&1\\0&0\end{bmatrix}}_{\text{shift }\mat{N}}.
\]
The shift part $\mat{N}=\begin{bmatrix}0&1\\0&0\end{bmatrix}$ satisfies $\mat{N}^2=\mat{0}$:
it is \emph{nilpotent}. Geometrically, $\mat{N}$ takes the second basis vector to the
first and the first to zero --- it pushes everything ``down the chain'' one step toward
the eigenvector and off the edge. The eigenvalue scales; the nilpotent slides. A Jordan
block is precisely \emph{scale-and-slide along a chain}, and that slide is the irreducible
core of non-diagonalizability.

\begin{figure}[H]
\centering
\begin{tikzpicture}[>=Stealth,scale=1.1,every node/.style={font=\small}]
  \node[draw=sourceblue, very thick, circle, fill=sourceblue!15, minimum size=1.0cm] (w) at (0,0) {$\vec{w}$};
  \node[draw=targetgreen, very thick, circle, fill=targetgreen!15, minimum size=1.0cm] (v) at (3.6,0) {$\vec{v}$};
  \node[draw=gray!70, thick, circle, fill=gray!15, minimum size=1.0cm] (z) at (7.2,0) {$\vec{0}$};

  \draw[->, very thick, structurepurple] (w) to[bend left=28] node[above=2pt, font=\footnotesize\bfseries] {$\mat{N}=\mat{A}-\lambda\mat{I}$} (v);
  \draw[->, very thick, structurepurple] (v) to[bend left=28] node[above=2pt, font=\footnotesize\bfseries] {$\mat{N}=\mat{A}-\lambda\mat{I}$} (z);

  \node[draw=sourceblue!40, rounded corners=2pt, fill=sourceblue!5, font=\footnotesize, align=center] at (0,-1.0)
    {Height 2\\[2pt]\scriptsize (generalized eigvec)};
  \node[draw=targetgreen!40, rounded corners=2pt, fill=targetgreen!5, font=\footnotesize, align=center] at (3.6,-1.0)
    {Height 1\\[2pt]\scriptsize (true eigenvector)};
  \node[draw=gray!40, rounded corners=2pt, fill=gray!5, font=\footnotesize, align=center] at (7.2,-1.0)
    {Height 0\\[2pt]\scriptsize (annihilated)};
\end{tikzpicture}
\caption{A Jordan chain of length $2$. The nilpotent part $\mat{N}=\mat{A}-\lambda\mat{I}$
walks the generalized eigenvector $\vec{w}$ down to the true eigenvector $\vec{v}$, then
down to $\vec{0}$. The block's size is the length of this staircase.}
\label{ch28:chain}
\end{figure}

\begin{aestheticbox}
Diagonalization shows that whenever an operator possesses a full eigenbasis, it acts as a set of
independent scalings along coordinate axes. The Jordan form shows that when
an eigenbasis is missing, the obstruction consists entirely of coupled chains
of generalized eigenvectors sliding into one another. Every square matrix over $\C$ decomposes into
a diagonal part and a commuting nilpotent part, where the nilpotent part decomposes into canonical
Jordan blocks. This reduces the classification of operators up to change of basis
to a catalog of eigenvalues together with integer partitions of block sizes.
\end{aestheticbox}

\begin{connectionbox}
The split $\mat{A}=\mat{D}+\mat{N}$ into a diagonalizable part and a commuting nilpotent
part ($\mat{D}\mat{N}=\mat{N}\mat{D}$) is the \emph{Jordan--Chevalley decomposition}, and
it echoes across mathematics. In Lie theory (the study of continuous geometric symmetries)
it separates an operator into its ``semisimple'' component (acting by pure scaling along
independent axes) and its ``unipotent'' component (an identity plus nilpotent part acting
by pure sliding). In differential equations it is why solutions to
$\vec{x}'=\mat{A}\vec{x}$ for a defective $\mat{A}$ pick up polynomial-times-exponential
terms like $t\,e^{\lambda t}$: the matrix exponential of a Jordan block is
\[
e^{t\mat{J}_2(\lambda)} = e^{\lambda t}\begin{bmatrix}1&t\\0&1\end{bmatrix},
\]
because $e^{t\mat{N}}=\mat{I}+t\mat{N}$ (the series stops since $\mat{N}^2=\mat{0}$). The
lone superdiagonal $1$ in the matrix becomes the lone $t$ in the solution. Resonance in an
undamped driven oscillator---where the response amplitude grows like $t$---is the physical
manifestation of a repeated eigenvalue with a nontrivial Jordan block.
\end{connectionbox}

\subsection{Nilpotent matrices are pure Jordan}

The nilpotent case deserves its own look, because it is the whole mystery with the
eigenvalue stripped away. A matrix $\mat{N}$ is \emph{nilpotent} if $\mat{N}^k=\mat{0}$
for some positive integer $k$---repeated applications crush every vector to zero. Then its only eigenvalue is $0$ (if $\mat{N}\vec{v}=\lambda\vec{v}$ then
$\vec{0}=\mat{N}^k\vec{v}=\lambda^k\vec{v}$, forcing $\lambda=0$), so its Jordan form is a
direct sum of blocks $\mat{J}_k(0)$ --- staircases of $1$'s with zeros on the diagonal.
Every defective behavior is, up to shifting by $\lambda\mat{I}$, a nilpotent shift. This
is why we kept saying ``the slide is the core'': the eigenvalue is just the diagonal
backdrop, and all the interesting structure lives in the nilpotent superdiagonal.

\begin{example}[Reading blocks from null-space growth]
\label{ch28:ex-nullgrowth}
Suppose a $4\times4$ nilpotent $\mat{N}$ has
\[
\dimn\nullsp\mat{N}=2,\quad \dimn\nullsp\mat{N}^2=3,\quad \dimn\nullsp\mat{N}^3=4,\quad \dimn\nullsp\mat{N}^4=4.
\]
So $d_1=2,d_2=3,d_3=4,d_4=4$ (with $d_0=0$). Number of blocks $=d_1=2$. Largest block
size $=3$ (where the null spaces stop growing). The blocks of size exactly $j$ number
$2d_j-d_{j-1}-d_{j+1}$: for $j=1$ that is $2(2)-0-3=1$; for $j=2$ it is $2(3)-2-4=0$; for
$j=3$ it is $2(4)-3-4=1$. So one block of size $1$ and one of size $3$: total $4$, two
blocks. The Jordan form is $\mat{J}_3(0)\oplus\mat{J}_1(0)$. We never computed an
eigenvector explicitly --- the dimensions of null spaces told us the whole structure.
\end{example}

\section{A Modern Application: Stability of Iterations}
\label{ch28:application}

\begin{applicationbox}
Many algorithms in machine learning and scientific computing are \emph{iterations}: a
state updates by $\vec{x}_{t+1}=\mat{A}\vec{x}_t$, so $\vec{x}_t=\mat{A}^t\vec{x}_0$.
Think of the power method for eigenvectors, a linear recurrent network, the propagation
of a gradient through the unrolled steps of a recurrent model, or a Markov chain mixing
to equilibrium. Whether the state blows up, decays, or drifts is governed entirely by the
powers $\mat{A}^t$ --- and the Jordan form is the cleanest way to understand $\mat{A}^t$.

If $\mat{A}$ is diagonalizable, $\mat{A}^t=\mat{P}\mat{D}^t\mat{P}\inv$ and each mode just
scales by $\lambda_i^{\,t}$: the spectral radius $\rho(\mat{A}) = \max_i|\lambda_i|$ (the largest eigenvalue magnitude, governing asymptotic exponential growth or decay) decides everything.
But for a defective $\mat{A}$ the powers of a Jordan block carry binomial coefficients:
\[
\mat{J}_2(\lambda)^t=\begin{bmatrix}\lambda^t & t\,\lambda^{t-1}\\[2pt]0&\lambda^t\end{bmatrix},
\qquad
\mat{J}_3(\lambda)^t=\begin{bmatrix}\lambda^t & t\lambda^{t-1} & \binom{t}{2}\lambda^{t-2}\\
0 & \lambda^t & t\lambda^{t-1}\\ 0&0&\lambda^t\end{bmatrix}.
\]
The off-diagonal entries grow \emph{polynomially in $t$} before the $\lambda^t$ takes
over. This is why an iteration with $|\lambda|$ \emph{exactly} $1$ and a nontrivial Jordan
block is unstable even though the spectral radius is $1$: the $t\,\lambda^{t-1}$ term
grows without bound. It is also why deep recurrent networks suffer ``exploding/vanishing
gradients'' that pure eigenvalue intuition underestimates --- near-defective weight
matrices inject transient polynomial growth into the backpropagated signal.
\end{applicationbox}

Those binomial coefficients come straight from the algebra. Expand
$\mat{J}_k(\lambda)=\lambda\mat{I}+\mat{N}$ with the binomial theorem, which is legitimate because
$\lambda\mat{I}$ and $\mat{N}$ commute:
\[
(\lambda\mat{I}+\mat{N})^t=\sum_{j=0}^{t}\binom{t}{j}\lambda^{t-j}\mat{N}^{j}.
\]
Since $\mat{N}^j$ is the matrix with $1$'s on the $j$-th superdiagonal and $\mat{N}^k=\mat{0}$,
only finitely many terms survive, and they land on successive superdiagonals with
coefficients $\binom{t}{j}\lambda^{t-j}$. That single line explains every entry of the
power formulas above.

\begin{recoverybox}
\textbf{Pause and think:} If every eigenvalue of $\mat{A}$ satisfies $|\lambda_i| < 1$, could a defective Jordan block prevent $\mat{A}^t$ from decaying to $\mat{0}$ as $t\to\infty$?
\emph{Reasoning:} No. Although the off-diagonal entries feature polynomial prefactors like $\binom{t}{j}\lambda^{t-j}$, exponential decay always conquers polynomial growth: for any fixed $j$ and any $|\lambda| < 1$, $\lim_{t\to\infty} t^j |\lambda|^t = 0$. The Jordan block causes a transient surge where $\|\mat{A}^t\|$ initially climbs, but the exponential decay eventually dominates and $\mat{A}^t \to \mat{0}$.
\end{recoverybox}

Figure~\ref{ch28:transient} makes the danger visible. With eigenvalue exactly on the unit
circle, $|\lambda|=1$, a diagonalizable matrix keeps every entry bounded forever, while a
single $2\times2$ Jordan block has an entry $t\,\lambda^{t-1}$ that grows without bound ---
slowly, linearly, but without bound. The spectral radius is the same in both cases; only
the block structure tells them apart.

\begin{figure}[H]
\centering
\begin{tikzpicture}
\begin{axis}[
  width=10cm,height=6cm,
  xlabel={iteration $t$},ylabel={size of off-diagonal entry},
  domain=0:30,samples=31,
  legend pos=north west,legend cell align=left,
  ymin=0,ymax=32,xmin=0,xmax=30,
  every axis plot/.append style={very thick},
  font=\small]
\addplot[vecblue]{1};
\addlegendentry{diagonalizable, $|\lambda|=1$ (bounded)}
\addplot[trapred]{x};
\addlegendentry{defective block, entry $t\,\lambda^{t-1}$, $|\lambda|=1$}
\end{axis}
\end{tikzpicture}
\caption{Why a defective eigenvalue on the unit circle is unstable. Both operators have
spectral radius exactly $1$, so eigenvalue intuition predicts ``neither grows nor shrinks.''
But the off-diagonal entry of a $2\times2$ Jordan block is $t\,\lambda^{t-1}$, which marches
upward linearly (red), while the diagonalizable case stays flat (blue). The transient is
invisible to the eigenvalues and visible only in the Jordan structure.}
\label{ch28:transient}
\end{figure}

\section{Putting It All Together}
\label{ch28:together}

Here is the recipe the whole chapter has been building, the one you can actually run on a
small matrix.

\begin{formalbox}
\textbf{Finding the Jordan form of $\mat{A}$.}
\begin{enumerate}
\item \textbf{Eigenvalues.} Factor the characteristic polynomial
  $\det(\mat{A}-t\mat{I})$. Each distinct root $\lambda$ with its multiplicity is an
  eigenvalue with its \emph{algebraic} multiplicity.
\item \textbf{Count blocks.} For each $\lambda$, compute
  $\dimn\nullsp(\mat{A}-\lambda\mat{I})$ --- the geometric multiplicity. That is the
  \emph{number of Jordan blocks} for $\lambda$.
\item \textbf{Size the blocks.} If you need the exact sizes, compute the growing
  nullities $d_j=\dimn\nullsp(\mat{A}-\lambda\mat{I})^j$ until they stop increasing, and
  read off block sizes from $2d_j-d_{j-1}-d_{j+1}$. (For small matrices the block sizes
  are usually forced once you know the algebraic and geometric multiplicities.)
\item \textbf{Build chains (if you want $\mat{P}$).} For each block, find a top
  generalized eigenvector $\vec{w}$ and apply $(\mat{A}-\lambda\mat{I})$ repeatedly to
  generate the chain down to a true eigenvector. Collect all chain vectors, bottom-up
  within each block, as the columns of $\mat{P}$.
\item \textbf{Result.} $\mat{A}=\mat{P}\mat{J}\mat{P}\inv$ with $\mat{J}$ the assembled
  blocks.
\end{enumerate}
\end{formalbox}

\begin{focusbox}
The one sentence to carry out of this chapter: \emph{similar matrices are one operator in
different bases, and the Jordan form is the simplest basis you can always reach --- a
diagonal of eigenvalues plus the minimum amount of shearing the operator genuinely
contains.} Diagonalizable is the case where that minimum is zero.
\end{focusbox}

\begin{remark}[A caution about the real and the numerical world]
\label{ch28:remark-caution}
Two cautions keep us grounded. First, Jordan form lives over $\C$: a real matrix may require complex
eigenvalues (for example, a planar rotation has eigenvalues $e^{\pm i\theta}$), so its Jordan form is complex
even though the original matrix is real. Second, the Jordan form is famously \emph{discontinuous}
and numerically unstable. The matrix $\begin{bmatrix}1&1\\0&1\end{bmatrix}$ has a $2\times2$
block, but $\begin{bmatrix}1&1\\0&1+\eps\end{bmatrix}$ for any $\eps\neq0$ is
diagonalizable with two distinct eigenvalues --- a tiny perturbation shatters the block.
This is why numerical libraries compute the \emph{Schur decomposition} (factoring $\mat{A} = \mat{Q}\mat{T}\mat{Q}^*$ using an orthonormal change of basis $\mat{Q}$ and an upper-triangular matrix $\mat{T}$) or the SVD instead, never the Jordan form. The Jordan form is an essential theoretical tool for analyzing powers and exponentials of matrices, even though floating-point software relies on the Schur decomposition or SVD instead.
\end{remark}

\begin{historybox}
The quest for canonical forms began in abstract algebra and Galois theory (the study of polynomial equations through the group of symmetries of their roots). In 1870, the French mathematician Camille Jordan published his masterwork, \emph{Trait\'e des substitutions et des \'equations alg\'ebriques}. Jordan was investigating symmetries of polynomial roots, which required classifying linear substitutions (what we now call linear transformations) over finite fields (number systems with finitely many elements where division is still possible) and over the complex numbers. When a matrix has distinct eigenvalues, ordinary diagonalization solves the classification problem completely: two matrices are similar if and only if they share the same spectrum (the complete set of eigenvalues). But for defective matrices, ordinary eigenvectors fail to span space. Jordan set out to find the simplest, most canonical form achievable for \emph{any} linear operator under a change of basis. Working with polynomials and determinants---decades before modern axiomatic vector spaces were formalized---Jordan proved that by augmenting eigenvectors with chains of generalized eigenvectors, any square matrix uncouples into a direct sum of Jordan blocks, cleanly separating the pure stretching of the eigenvalues from the irreducible nilpotent shearing.

A heated 19th-century priority controversy soon followed. Two years earlier, in 1868, the German mathematician Karl Weierstrass had published a classification of pairs of bilinear forms using polynomial invariants he called ``elementary divisors'' (determinantal divisors of matrix polynomials). Weierstrass and his supporters argued that Jordan had merely translated Weierstrass's theory into the language of substitutions. Jordan countered that his block-matrix form was explicit, constructive, and provided the geometric basis that Weierstrass's invariant polynomials obscured. History ultimately credited both: Weierstrass provided the coordinate-free polynomial invariants, while Jordan gave the unforgettable visual matrix.

A century later, a second controversy erupted at the intersection of pure algebra and computational physics. Pure mathematicians celebrated the Jordan form as the definitive structural classification of linear maps. But in the 1960s, British numerical analyst James H. Wilkinson demonstrated that the Jordan canonical form is fundamentally ill-posed and uncomputable in floating-point arithmetic. The issue is severe discontinuity: consider the matrix $\begin{bmatrix}0&1\\\eps&0\end{bmatrix}$. When $\eps = 0$, the matrix has eigenvalue $0$ with only one eigenvector---a defective $2\times 2$ Jordan block. But if $\eps = 10^{-16}$, the magnitude of a single roundoff error in double-precision arithmetic, the eigenvalues split into $+10^{-8}$ and $-10^{-8}$, and the Jordan block shatters into a diagonalizable matrix. Because finite precision cannot distinguish whether two floating-point numbers are truly equal or merely close, Jordan form cannot be determined reliably by an algorithm. In 1909, Issai Schur had already discovered the numerical antidote: the \emph{Schur decomposition} ($\mat{A} = \mat{Q}\mat{T}\mat{Q}^*$, where $\mat{Q}$ is unitary, preserving lengths without distortion, and $\mat{T}$ is upper triangular). Numerical software computes Schur forms, never Jordan forms.

What Jordan form unlocked was the complete classification of linear operators up to similarity, alongside an exact explanation for why defective systems produce secular terms (polynomial growth factors like $t^k$ that multiply exponentials over time in differential equations). What remains open? The problem of \emph{simultaneous similarity}: finding a canonical form for a pair of non-commuting matrices $(\mat{A}, \mat{B})$ under a shared basis change $\mat{P}\inv\mat{A}\mat{P}, \mat{P}\inv\mat{B}\mat{P}$. In modern representation theory, this is proven to be a ``wild'' problem---an intractable classification task equivalent to classifying representations of all finite-dimensional algebras. In applied mathematics, finding robust, stable surrogates for defective matrix dynamics (such as the pseudospectrum, which tracks how eigenvalues behave under small perturbations) remains an active and vital frontier.
\end{historybox}

\subsection*{Looking forward: from similarity to the SVD}

Throughout this chapter, similarity forced us to play by a strict rule: we had to use the \emph{same} change of basis $\mat{P}$ for both the input space and the output space ($\mat{B} = \mat{P}\inv\mat{A}\mat{P}$). That rigid demand is the single reason why defective matrices cannot be diagonalized, and why Jordan form requires potentially fragile, non-orthogonal bases.

In Chapter~29, we relax that constraint. By allowing \emph{two different} orthonormal bases---one for inputs ($\mat{V}$) and one for outputs ($\mat{U}$)---the Singular Value Decomposition ($\mat{A} = \mat{U}\mat{\Sigma}\mat{V}\T$) achieves a diagonal representation for \emph{every} matrix, whether square, rectangular, defective, or well-behaved, with complete numerical stability. While similarity is restricted to single-basis changes on square matrices, the SVD provides an orthogonal factorization for any linear map.

\section{Exercises}
\label{ch28:exercises}

\begin{exercise}[Warmup: invariants]
\label{ch28:ex-invariants}
Let $\mat{A}=\begin{bmatrix}4&2\\1&3\end{bmatrix}$ and
$\mat{P}=\begin{bmatrix}1&1\\0&1\end{bmatrix}$. Compute $\mat{B}=\mat{P}\inv\mat{A}\mat{P}$.
Verify by direct calculation that $\tr\mat{B}=\tr\mat{A}$, $\det\mat{B}=\det\mat{A}$, and
that $\mat{A}$ and $\mat{B}$ have the same characteristic polynomial.
\end{exercise}

\begin{exercise}[Geometric vs.\ algebraic multiplicity]
\label{ch28:ex-mult}
For each matrix, find the eigenvalues and, for each eigenvalue, its algebraic and
geometric multiplicity. Which are diagonalizable?
\[
\text{(a)}\ \begin{bmatrix}2&0\\0&2\end{bmatrix}\quad
\text{(b)}\ \begin{bmatrix}2&1\\0&2\end{bmatrix}\quad
\text{(c)}\ \begin{bmatrix}3&0&0\\0&3&1\\0&0&3\end{bmatrix}.
\]
\end{exercise}

\begin{exercise}[A chain by hand]
\label{ch28:ex-chain}
Let $\mat{A}=\begin{bmatrix}6&1\\-1&4\end{bmatrix}$. Show its only eigenvalue is $5$ with
geometric multiplicity $1$. Find the eigenvector $\vec{v}$, then a generalized
eigenvector $\vec{w}$ solving $(\mat{A}-5\mat{I})\vec{w}=\vec{v}$. Form
$\mat{P}=[\vec{v}\ \vec{w}]$ and verify $\mat{P}\inv\mat{A}\mat{P}=\begin{bmatrix}5&1\\0&5\end{bmatrix}$.
\end{exercise}

\begin{exercise}[Powers of a block]
\label{ch28:ex-powers}
Using $\mat{J}_2(\lambda)=\lambda\mat{I}+\mat{N}$ and $\mat{N}^2=\mat{0}$, prove by the
binomial theorem that $\mat{J}_2(\lambda)^t=\begin{bmatrix}\lambda^t & t\lambda^{t-1}\\
0&\lambda^t\end{bmatrix}$. Then evaluate $\mat{J}_2(1)^{10}$ and $\mat{J}_2(\tfrac12)^{10}$
and comment on which iteration is stable.
\end{exercise}

\begin{exercise}[Same eigenvalues, not similar]
\label{ch28:ex-notsimilar}
Explain why $\begin{bmatrix}2&0&0\\0&2&0\\0&0&2\end{bmatrix}$,
$\begin{bmatrix}2&1&0\\0&2&0\\0&0&2\end{bmatrix}$, and
$\begin{bmatrix}2&1&0\\0&2&1\\0&0&2\end{bmatrix}$ all share the characteristic polynomial
$(t-2)^3$ yet no two of them are similar. Identify the geometric multiplicity of $2$ in
each, and state each matrix's Jordan form.
\end{exercise}

\begin{exercise}[Reading blocks from nullities]
\label{ch28:ex-nullities}
A $5\times5$ matrix has a single eigenvalue $\lambda$ with
$d_1=\dimn\nullsp(\mat{A}-\lambda\mat{I})=3$, $d_2=4$, $d_3=5$, and $d_4=5$. How many
Jordan blocks are there? What is the size of the largest block? Determine the complete
list of block sizes and write the Jordan form. (Hint: use
$2d_j-d_{j-1}-d_{j+1}$ with $d_0=0$.)
\end{exercise}

\begin{exercise}[Nilpotent practice]
\label{ch28:ex-nilpotent}
Let $\mat{N}=\begin{bmatrix}0&1&2\\0&0&3\\0&0&0\end{bmatrix}$. Confirm $\mat{N}^3=\mat{0}$
but $\mat{N}^2\neq\mat{0}$. What is the only eigenvalue? Using the height of the chain,
argue that the Jordan form is a single block $\mat{J}_3(0)$.
\end{exercise}

\begin{exercise}[Differential equations connection]
\label{ch28:ex-ode}
For $\mat{A}=\begin{bmatrix}\lambda&1\\0&\lambda\end{bmatrix}$, use
$e^{t\mat{A}}=e^{\lambda t}(\mat{I}+t\mat{N})$ to write the general solution of
$\vec{x}'=\mat{A}\vec{x}$. Explain why a defective eigenvalue produces a
$t\,e^{\lambda t}$ term, and connect this to resonance in a forced oscillator.
\end{exercise}

\begin{exercise}[Proof: trace and similarity]
\label{ch28:ex-trace-proof}
Prove that $\tr(\mat{X}\mat{Y})=\tr(\mat{Y}\mat{X})$ for any conformable matrices, and use
it to give a one-line proof that similar matrices have equal trace. Then explain why the
converse fails: give two matrices with equal trace that are not similar.
\end{exercise}

\begin{exercise}[Challenge: the minimal polynomial]
\label{ch28:ex-minpoly}
The \emph{minimal polynomial} of $\mat{A}$ is the smallest-degree monic polynomial $m(t)$
with $m(\mat{A})=\mat{0}$. Argue that for each eigenvalue $\lambda$, the exponent of
$(t-\lambda)$ in $m(t)$ equals the size of the \emph{largest} Jordan block for $\lambda$.
Use this to show $\mat{A}$ is diagonalizable if and only if its minimal polynomial has no
repeated roots. (Hint: think about what power of $\mat{N}=\mat{J}_k(\lambda)-\lambda\mat{I}$
first becomes zero.)
\end{exercise}

\begin{exercise}[Challenge: classify $3\times3$ up to similarity]
\label{ch28:ex-classify}
Suppose a $3\times3$ complex matrix has characteristic polynomial $(t-1)^2(t-2)$. List
all possible Jordan forms (there are exactly two), describe how to tell them apart using
the geometric multiplicity of the eigenvalue $1$, and state for each whether the matrix is
diagonalizable.
\end{exercise}
